\section{引言}

课程论文通常用于展示学生对某门研究生课程核心概念、模型方法和实证分析流程的理解。本文以应用统计专业常见的时间序列预测问题为例，展示如何使用本模板完成封面、目录、正文、图表、公式、参考文献和附代码。

写作时建议保持“问题背景--模型方法--数据分析--结论讨论”的结构。若课程要求更偏理论推导，可以把第二部分改写为定理、证明和模拟实验；若课程要求更偏应用，可以增加数据来源、变量说明、模型诊断和稳健性检验。

本文示例围绕居民消费价格指数 \(y_t\) 的动态建模展开。设观测序列为 \(\{y_t:t=1,\ldots,T\}\)，研究目标是在给定历史信息集 \(\mathcal{F}_{t-1}\) 的条件下，对未来一期或多期指标进行预测：
\begin{equation}
  \hat{y}_{t+h\mid t}=\E(y_{t+h}\mid \mathcal{F}_{t}),\qquad h=1,2,\ldots,H.
  \label{eq:forecast-target}
\end{equation}

如 \cref{eq:forecast-target} 所示，统计建模不仅要给出点预测，还应尽可能给出区间预测、误差分解和模型不确定性解释。
